% !TEX program = xelatex
\documentclass[a4paper,11pt,fontset=none]{ctexart}

\newcommand{\reporttitle}{华润微 vs 士兰微：功率半导体双龙头深度对比}
\newcommand{\reportsubtitle}{688396 / 600460 | 功率IDM | AI+新能源+SiC共振}
\newcommand{\reportkicker}{ASTOCK RESEARCH NOTE}
\newcommand{\reportscope}{CHINA A-SHARES | POWER SEMICONDUCTOR}
\newcommand{\reportdate}{2026-07-06}
\newcommand{\reportdatacutoff}{2026-07-06 (Web search aggregated)}
\newcommand{\reporttype}{Multi-agent comparative research note}
\newcommand{\reportauthor}{AStock Agent Team}
\newcommand{\reporthouseview}{%
华润微与士兰微均处于功率半导体行业景气上行周期中，2026Q1业绩均大幅改善（华润微净利+297\%，士兰微净利+41\%）。华润微定位``央企IDM全产业链龙头''，MOSFET国内第一，AI服务器电源为核心增量；士兰微定位``民企高端功率IDM''，白电IPM全球第一，SiC车规为第二曲线。两家公司各具差异化优势，估值均处于较高分位。华润微当前PE TTM$\sim$127x（市值1150亿），士兰微PE TTM$\sim$169x（市值776亿），均需业绩持续高增验证。值得中长期跟踪，但短期追高风险较大。%
}
\newcommand{\reportquality}{Web search aggregated; local data pipeline unavailable. Quality tier: partial source / snapshot degraded.}
\newcommand{\reportdisclaimer}{本报告不构成任何投资建议。}
% Reusable IB-style preamble for XeLaTeX equity research reports
% Usage: % Reusable IB-style preamble for XeLaTeX equity research reports
% Usage: % Reusable IB-style preamble for XeLaTeX equity research reports
% Usage: \input{preamble} after \documentclass[a4paper,11pt,openany,fontset=none]{ctexrep}

% ============ 页面布局 ============
\usepackage[top=2cm, bottom=2cm, left=2cm, right=2cm]{geometry}
\setlength{\parskip}{0.3em}
\setlength{\parindent}{2em}
\renewcommand{\baselinestretch}{1.15}

% ============ 字体（macOS） ============
\setCJKmainfont[BoldFont={Heiti SC}]{STSong}
\setCJKsansfont{Heiti SC}
\setCJKmonofont{Heiti SC}
\setmainfont{Times New Roman}

% ============ 颜色（IB Navy/Grey palette） ============
\usepackage{xcolor}
\definecolor{navy}{HTML}{003366}
\definecolor{steelgrey}{HTML}{4A5568}
\definecolor{lightgrey}{HTML}{F7FAFC}
\definecolor{riskred}{HTML}{C53030}
\definecolor{riskamber}{HTML}{D69E2E}
\definecolor{riskgreen}{HTML}{38A169}
\definecolor{nvgreen}{HTML}{76B900}

% ============ 表格 ============
\usepackage{booktabs}
\usepackage{longtable}
\usepackage{tabularx}
\usepackage{multirow}
\usepackage{array}
\newcolumntype{Y}{>{\centering\arraybackslash}X}
\newcolumntype{L}[1]{>{\raggedright\arraybackslash}p{#1}}
\newcolumntype{C}[1]{>{\centering\arraybackslash}p{#1}}
\newcolumntype{R}[1]{>{\raggedleft\arraybackslash}p{#1}}

% ============ 图形 ============
\usepackage{tikz}
\usetikzlibrary{shapes,arrows.meta,positioning,calc,backgrounds,fit}
\usepackage{pgfplots}
\usepgfplotslibrary{polar}
\pgfplotsset{compat=1.18}
\usepackage[edges]{forest}

% ============ 页眉页脚（报告标题和日期需在main.tex中覆盖） ============
\usepackage{fancyhdr}
\pagestyle{fancy}
\fancyhf{}
\fancyhead[L]{\small\color{steelgrey} \reporttitle}
\fancyhead[R]{\small\color{steelgrey} \reportdate}
\fancyfoot[C]{\small\color{steelgrey} \thepage}
\fancyfoot[R]{\tiny\color{steelgrey} 本报告不构成任何投资建议}
\renewcommand{\headrulewidth}{0.4pt}
\renewcommand{\footrulewidth}{0.2pt}

% ============ 章节格式 ============
\usepackage{titlesec}
\titleformat{\chapter}[display]
  {\normalfont\huge\bfseries\color{navy}}
  {\chaptertitlename\ \thechapter}{20pt}{\Huge}
\titleformat{\section}
  {\normalfont\Large\bfseries\color{navy}}
  {\thesection}{1em}{}
\titleformat{\subsection}
  {\normalfont\large\bfseries\color{steelgrey}}
  {\thesubsection}{1em}{}

% ============ 超链接 ============
\usepackage{hyperref}
\hypersetup{
  colorlinks=true,
  linkcolor=navy,
  urlcolor=navy,
  citecolor=navy,
}

% ============ 其他工具 ============
\usepackage{enumitem}
\setlist{nosep, leftmargin=2em}
\usepackage{fontawesome5}
\usepackage{amssymb}
\usepackage{pifont}
\usepackage{tcolorbox}
\tcbuselibrary{skins,breakable}
\usepackage{caption}
\captionsetup{font=small, skip=4pt}
\usepackage{float}
\setlength{\textfloatsep}{8pt plus 2pt minus 2pt}
\setlength{\floatsep}{6pt plus 2pt minus 2pt}
\setlength{\intextsep}{6pt plus 2pt minus 2pt}
\setlength{\abovecaptionskip}{4pt}
\setlength{\belowcaptionskip}{2pt}

% ============ 自定义命令 ============
\newcommand{\rmark}{\textcolor{riskred}{\ding{108}}}
\newcommand{\amark}{\textcolor{riskamber}{\ding{108}}}
\newcommand{\gmark}{\textcolor{riskgreen}{\ding{108}}}
\newcommand{\starfive}{$\bigstar\bigstar\bigstar\bigstar\bigstar$}
\newcommand{\starfour}{$\bigstar\bigstar\bigstar\bigstar$}
\newcommand{\starthree}{$\bigstar\bigstar\bigstar$}

% 信息框
\newtcolorbox{keyinsight}[1][]{
  colback=lightgrey, colframe=navy,
  fonttitle=\bfseries, title={#1}, breakable
}
\newtcolorbox{riskbox}[1][]{
  colback=red!5, colframe=riskred,
  fonttitle=\bfseries, title={#1}, breakable
}
 after \documentclass[a4paper,11pt,openany,fontset=none]{ctexrep}

% ============ 页面布局 ============
\usepackage[top=2cm, bottom=2cm, left=2cm, right=2cm]{geometry}
\setlength{\parskip}{0.3em}
\setlength{\parindent}{2em}
\renewcommand{\baselinestretch}{1.15}

% ============ 字体（macOS） ============
\setCJKmainfont[BoldFont={Heiti SC}]{STSong}
\setCJKsansfont{Heiti SC}
\setCJKmonofont{Heiti SC}
\setmainfont{Times New Roman}

% ============ 颜色（IB Navy/Grey palette） ============
\usepackage{xcolor}
\definecolor{navy}{HTML}{003366}
\definecolor{steelgrey}{HTML}{4A5568}
\definecolor{lightgrey}{HTML}{F7FAFC}
\definecolor{riskred}{HTML}{C53030}
\definecolor{riskamber}{HTML}{D69E2E}
\definecolor{riskgreen}{HTML}{38A169}
\definecolor{nvgreen}{HTML}{76B900}

% ============ 表格 ============
\usepackage{booktabs}
\usepackage{longtable}
\usepackage{tabularx}
\usepackage{multirow}
\usepackage{array}
\newcolumntype{Y}{>{\centering\arraybackslash}X}
\newcolumntype{L}[1]{>{\raggedright\arraybackslash}p{#1}}
\newcolumntype{C}[1]{>{\centering\arraybackslash}p{#1}}
\newcolumntype{R}[1]{>{\raggedleft\arraybackslash}p{#1}}

% ============ 图形 ============
\usepackage{tikz}
\usetikzlibrary{shapes,arrows.meta,positioning,calc,backgrounds,fit}
\usepackage{pgfplots}
\usepgfplotslibrary{polar}
\pgfplotsset{compat=1.18}
\usepackage[edges]{forest}

% ============ 页眉页脚（报告标题和日期需在main.tex中覆盖） ============
\usepackage{fancyhdr}
\pagestyle{fancy}
\fancyhf{}
\fancyhead[L]{\small\color{steelgrey} \reporttitle}
\fancyhead[R]{\small\color{steelgrey} \reportdate}
\fancyfoot[C]{\small\color{steelgrey} \thepage}
\fancyfoot[R]{\tiny\color{steelgrey} 本报告不构成任何投资建议}
\renewcommand{\headrulewidth}{0.4pt}
\renewcommand{\footrulewidth}{0.2pt}

% ============ 章节格式 ============
\usepackage{titlesec}
\titleformat{\chapter}[display]
  {\normalfont\huge\bfseries\color{navy}}
  {\chaptertitlename\ \thechapter}{20pt}{\Huge}
\titleformat{\section}
  {\normalfont\Large\bfseries\color{navy}}
  {\thesection}{1em}{}
\titleformat{\subsection}
  {\normalfont\large\bfseries\color{steelgrey}}
  {\thesubsection}{1em}{}

% ============ 超链接 ============
\usepackage{hyperref}
\hypersetup{
  colorlinks=true,
  linkcolor=navy,
  urlcolor=navy,
  citecolor=navy,
}

% ============ 其他工具 ============
\usepackage{enumitem}
\setlist{nosep, leftmargin=2em}
\usepackage{fontawesome5}
\usepackage{amssymb}
\usepackage{pifont}
\usepackage{tcolorbox}
\tcbuselibrary{skins,breakable}
\usepackage{caption}
\captionsetup{font=small, skip=4pt}
\usepackage{float}
\setlength{\textfloatsep}{8pt plus 2pt minus 2pt}
\setlength{\floatsep}{6pt plus 2pt minus 2pt}
\setlength{\intextsep}{6pt plus 2pt minus 2pt}
\setlength{\abovecaptionskip}{4pt}
\setlength{\belowcaptionskip}{2pt}

% ============ 自定义命令 ============
\newcommand{\rmark}{\textcolor{riskred}{\ding{108}}}
\newcommand{\amark}{\textcolor{riskamber}{\ding{108}}}
\newcommand{\gmark}{\textcolor{riskgreen}{\ding{108}}}
\newcommand{\starfive}{$\bigstar\bigstar\bigstar\bigstar\bigstar$}
\newcommand{\starfour}{$\bigstar\bigstar\bigstar\bigstar$}
\newcommand{\starthree}{$\bigstar\bigstar\bigstar$}

% 信息框
\newtcolorbox{keyinsight}[1][]{
  colback=lightgrey, colframe=navy,
  fonttitle=\bfseries, title={#1}, breakable
}
\newtcolorbox{riskbox}[1][]{
  colback=red!5, colframe=riskred,
  fonttitle=\bfseries, title={#1}, breakable
}
 after \documentclass[a4paper,11pt,openany,fontset=none]{ctexrep}

% ============ 页面布局 ============
\usepackage[top=2cm, bottom=2cm, left=2cm, right=2cm]{geometry}
\setlength{\parskip}{0.3em}
\setlength{\parindent}{2em}
\renewcommand{\baselinestretch}{1.15}

% ============ 字体（macOS） ============
\setCJKmainfont[BoldFont={Heiti SC}]{STSong}
\setCJKsansfont{Heiti SC}
\setCJKmonofont{Heiti SC}
\setmainfont{Times New Roman}

% ============ 颜色（IB Navy/Grey palette） ============
\usepackage{xcolor}
\definecolor{navy}{HTML}{003366}
\definecolor{steelgrey}{HTML}{4A5568}
\definecolor{lightgrey}{HTML}{F7FAFC}
\definecolor{riskred}{HTML}{C53030}
\definecolor{riskamber}{HTML}{D69E2E}
\definecolor{riskgreen}{HTML}{38A169}
\definecolor{nvgreen}{HTML}{76B900}

% ============ 表格 ============
\usepackage{booktabs}
\usepackage{longtable}
\usepackage{tabularx}
\usepackage{multirow}
\usepackage{array}
\newcolumntype{Y}{>{\centering\arraybackslash}X}
\newcolumntype{L}[1]{>{\raggedright\arraybackslash}p{#1}}
\newcolumntype{C}[1]{>{\centering\arraybackslash}p{#1}}
\newcolumntype{R}[1]{>{\raggedleft\arraybackslash}p{#1}}

% ============ 图形 ============
\usepackage{tikz}
\usetikzlibrary{shapes,arrows.meta,positioning,calc,backgrounds,fit}
\usepackage{pgfplots}
\usepgfplotslibrary{polar}
\pgfplotsset{compat=1.18}
\usepackage[edges]{forest}

% ============ 页眉页脚（报告标题和日期需在main.tex中覆盖） ============
\usepackage{fancyhdr}
\pagestyle{fancy}
\fancyhf{}
\fancyhead[L]{\small\color{steelgrey} \reporttitle}
\fancyhead[R]{\small\color{steelgrey} \reportdate}
\fancyfoot[C]{\small\color{steelgrey} \thepage}
\fancyfoot[R]{\tiny\color{steelgrey} 本报告不构成任何投资建议}
\renewcommand{\headrulewidth}{0.4pt}
\renewcommand{\footrulewidth}{0.2pt}

% ============ 章节格式 ============
\usepackage{titlesec}
\titleformat{\chapter}[display]
  {\normalfont\huge\bfseries\color{navy}}
  {\chaptertitlename\ \thechapter}{20pt}{\Huge}
\titleformat{\section}
  {\normalfont\Large\bfseries\color{navy}}
  {\thesection}{1em}{}
\titleformat{\subsection}
  {\normalfont\large\bfseries\color{steelgrey}}
  {\thesubsection}{1em}{}

% ============ 超链接 ============
\usepackage{hyperref}
\hypersetup{
  colorlinks=true,
  linkcolor=navy,
  urlcolor=navy,
  citecolor=navy,
}

% ============ 其他工具 ============
\usepackage{enumitem}
\setlist{nosep, leftmargin=2em}
\usepackage{fontawesome5}
\usepackage{amssymb}
\usepackage{pifont}
\usepackage{tcolorbox}
\tcbuselibrary{skins,breakable}
\usepackage{caption}
\captionsetup{font=small, skip=4pt}
\usepackage{float}
\setlength{\textfloatsep}{8pt plus 2pt minus 2pt}
\setlength{\floatsep}{6pt plus 2pt minus 2pt}
\setlength{\intextsep}{6pt plus 2pt minus 2pt}
\setlength{\abovecaptionskip}{4pt}
\setlength{\belowcaptionskip}{2pt}

% ============ 自定义命令 ============
\newcommand{\rmark}{\textcolor{riskred}{\ding{108}}}
\newcommand{\amark}{\textcolor{riskamber}{\ding{108}}}
\newcommand{\gmark}{\textcolor{riskgreen}{\ding{108}}}
\newcommand{\starfive}{$\bigstar\bigstar\bigstar\bigstar\bigstar$}
\newcommand{\starfour}{$\bigstar\bigstar\bigstar\bigstar$}
\newcommand{\starthree}{$\bigstar\bigstar\bigstar$}

% 信息框
\newtcolorbox{keyinsight}[1][]{
  colback=lightgrey, colframe=navy,
  fonttitle=\bfseries, title={#1}, breakable
}
\newtcolorbox{riskbox}[1][]{
  colback=red!5, colframe=riskred,
  fonttitle=\bfseries, title={#1}, breakable
}


\begin{document}

\begin{center}
{\sffamily\small\bfseries\color{steelgrey} \reportkicker}\\[3pt]
{\LARGE\bfseries\color{navy} \reporttitle}\\[5pt]
{\large\color{steelgrey} \reportsubtitle}\\[4pt]
{\small\color{mutedgrey} \reportscope \quad|\quad \reportdate \quad|\quad \reportdatacutoff}
\end{center}
\vspace{0.7em}

\begin{dashboardbox}[Decision Snapshot]
\begin{houseviewbox}[House View]
\reporthouseview
\end{houseviewbox}
\begin{sourcequalitybox}[Data Quality]
\reportquality
\end{sourcequalitybox}
\end{dashboardbox}

%% ================================================================
%% PART I: 华润微
%% ================================================================
\newpage
\section*{\color{navy}\large Part I: 华润微（688396.SH）}
\addcontentsline{toc}{section}{Part I: 华润微}

\section{公司概况与核心定位}

\textbf{华润微}（688396.SH）是中国功率半导体IDM绝对龙头，央企背景（华润集团），是国内少数具备芯片设计、掩模制造、晶圆制造、封装测试全产业链一体化运营能力的半导体企业。

\textbf{核心业务}：
\begin{itemize}
  \item \textbf{功率器件}（核心，约60\%+）：MOSFET（国内市占率>20\%，第一）、IGBT、晶闸管。
  \item \textbf{智能传感器}：光电、压力、MCU等。
  \item \textbf{制造代工}：6/8/12英寸晶圆代工（重庆、无锡、深圳）。
  \item \textbf{第三代半导体}：SiC MOSFET/模块、GaN产品。
\end{itemize}

\textbf{核心壁垒}：
\begin{itemize}
  \item 国内功率半导体IDM唯一龙头，MOSFET本土市占率第一（>20\%）。
  \item 央企背景+全产业链覆盖，客户认证壁垒高。
  \item AI服务器电源全链路覆盖（DrMOS + SGT MOS + SiC + 控制IC + 模块一体化方案）。
  \item 6/8/12英寸全系列产线，重庆12英寸已满产，深圳加速扩产。
\end{itemize}

\section{华润微业绩分析}

\begin{exhibitbox}[Exhibit 1: 华润微业绩拐点]
\begin{tabularx}{\textwidth}{L{3.2cm}XXXX}
\toprule
\textbf{指标} & \textbf{2024A} & \textbf{2025A} & \textbf{2026Q1} & \textbf{2026E(机构)} \\
\midrule
营收（亿元） & 86.7 & $\sim$100 & 28.57（+21\%） & 120--130 \\
归母净利（亿元） & 6.6 & 9.59 & 3.30（+297\%） & 12--13.5 \\
毛利率 & 23\% & $\sim$27\% & 25.48\% & 32\%+ \\
经营现金流（亿） & --- & --- & 6.27（+132\%） & --- \\
\bottomrule
\end{tabularx}
\sourcenote{数据来源：公司公告、华鑫证券（2026.06.21）、华源证券}
\end{exhibitbox}

\textbf{关键判断}：
\begin{itemize}
  \item 2026Q1净利暴增297\%，确认功率半导体景气回暖+提价+产能利用率提升三重共振。
  \item AI服务器电源订单占比提升至约30\%，SGT MOS/超级MOS已批量供货头部OEM。
  \item 功率半导体涨价：订单能见度最高9个月，第二轮涨价谈判已启动。
  \item 华源证券预测2025--2027E归母净利9.59/12.05/16.50亿元（+25.8\%/+25.6\%/+36.9\%）。
  \item 高盛认为AI电源方面将有更多贡献，光模块电源驱动模块预计2026H2大规模生产。
\end{itemize}

\section{华润微行业地位与竞争格局}

\begin{tabularx}{\textwidth}{L{3cm}L{3cm}X}
\toprule
\textbf{细分领域} & \textbf{市场地位} & \textbf{竞争对手} \\
\midrule
MOSFET & 国内第一（>20\%） & 英飞凌、安森美、士兰微 \\
功率IDM综合 & 国内唯一全产线龙头 & 士兰微（部分重叠） \\
AI服务器电源MOS & 核心供应商 & TI、英飞凌 \\
SiC MOSFET & 国内第一梯队 & 三安光电、士兰微 \\
车规功率 & 全品类认证 & 英飞凌、士兰微、斯达半导 \\
晶闸管/工控 & 国内绝对龙头 & --- \\
\bottomrule
\end{tabularx}

\textbf{差异化}：MOSFET国产化率仅$\sim$15\%，IGBT/SiC约35\%，替代空间巨大。华润微作为央企IDM具备客户信任度和规模优势。

\section{华润微增长驱动}

\begin{enumerate}
  \item \textbf{AI服务器电源}：提供PSU/HVDC/SST/IBC/POL全解决方案，DrMOS+控制器一体化，绑定头部服务器OEM。
  \item \textbf{功率半导体涨价}：行业进入上行周期，订单能见度9个月，第二轮涨价启动。
  \item \textbf{车规SiC}：SiC MOSFET覆盖650V--2300V，已批量上车，3300V规划中。
  \item \textbf{产能释放}：重庆12英寸满产，深圳产线加速扩产，规模效应持续释放。
  \item \textbf{国产替代}：功率半导体国产化率低位，央企IDM为国产替代首选。
  \item \textbf{光模块电源}：与领先光模块厂商合作，2026H2大规模生产。
\end{enumerate}

\section{华润微估值与市场定价}

\begin{tabularx}{\textwidth}{L{3cm}X}
\toprule
\textbf{Item} & \textbf{Value} \\
\midrule
最新股价 & 86.03元（2026-07-06），近期高点88.33元 \\
总市值 & $\sim$1152亿元 \\
PE TTM & $\sim$127x \\
PE（2026E） & 85--96x（基于12--13.5亿净利） \\
PE（2027E） & 70x（基于16.5亿净利） \\
PB & 4.26x \\
近期走势 & 整体强势上涨趋势中，7/6日+2.91\% \\
\bottomrule
\end{tabularx}

%% ================================================================
%% PART II: 士兰微
%% ================================================================
\newpage
\section*{\color{navy}\large Part II: 士兰微（600460.SH）}
\addcontentsline{toc}{section}{Part II: 士兰微}

\section{公司概况与核心定位}

\textbf{士兰微}（600460.SH）是国内IDM模式功率半导体龙头（民企），全球功率半导体市场排名第六（市占率3.3\%），深度绑定新能源车、光伏储能、白色家电三大赛道。

\textbf{核心业务}：
\begin{itemize}
  \item \textbf{功率器件与模块}（核心，$\sim$70\%）：IGBT单管/模块、超结MOSFET、IPM智能功率模块。
  \item \textbf{集成电路}（$\sim$20\%）：模拟IC、MCU、LED驱动。
  \item \textbf{第三代半导体}（快速增长）：SiC MOSFET（车规800V平台）、GaN器件。
\end{itemize}

\textbf{核心壁垒}：
\begin{itemize}
  \item 白电变频IPM全球第一（国内白电市占率>40\%），光伏MOS国内第一梯队。
  \item 全球功率半导体排名第六，国内市占率第一。
  \item 8英寸成熟产线满产盈利+12英寸高端模拟产线（200亿投资）+8英寸SiC产线三线并进。
  \item 车规全认证（IATF16949、ISO26262 ASILD、TISAX），IGBT/SiC深度绑定主流车企。
\end{itemize}

\section{士兰微业绩分析}

\begin{exhibitbox}[Exhibit 2: 士兰微业绩修复]
\begin{tabularx}{\textwidth}{L{3.2cm}XXXX}
\toprule
\textbf{指标} & \textbf{2024A} & \textbf{2025A} & \textbf{2026Q1} & \textbf{2026E(机构)} \\
\midrule
营收（亿元） & 112.2 & 130.52（+16\%） & 35.19（+17\%） & 155 \\
归母净利（亿元） & 0.90 & 3.99（+345\%） & 2.09（+41\%） & 6--8 \\
扣非净利（亿元） & --- & --- & 1.98（+45\%） & --- \\
毛利率 & --- & 逐季修复 & 持续改善 & --- \\
\bottomrule
\end{tabularx}
\sourcenote{数据来源：公司公告、华泰证券（2026.04.28）、国信证券、凯基证券}
\end{exhibitbox}

\textbf{关键判断}：
\begin{itemize}
  \item 2025年归母净利同比+345\%（从底部大幅反弹），2026Q1延续高增。
  \item Q1营收35.19亿（+17\% YoY，+33\% QoQ），环比改善显著。
  \item 8/12英寸硅基功率与模拟产线全年满载运行，12英寸产出同比+19\%。
  \item 机构预测2026--2028年营收155/195/245亿（+19\%/+26\%/+26\%）。
  \item 凯基证券给予目标价55元（当前46.62元，潜在上行空间18\%）。
\end{itemize}

\section{士兰微行业地位与竞争格局}

\begin{tabularx}{\textwidth}{L{3cm}L{3cm}X}
\toprule
\textbf{细分领域} & \textbf{市场地位} & \textbf{竞争对手} \\
\midrule
白电IPM & 全球第四/国内第一（>40\%） & 三菱、富士电机 \\
功率半导体综合 & 全球第六（3.3\%） & 英飞凌、安森美、华润微 \\
车规IGBT & 国内第一梯队 & 斯达半导、华润微、英飞凌 \\
SiC车规 & 全球第六 & Wolfspeed、ST、华润微 \\
光伏储能功率 & 国内第一梯队 & 华润微、阳光电源自供 \\
中小功率MOS & 国内龙头 & 华润微、新洁能 \\
\bottomrule
\end{tabularx}

\section{士兰微增长驱动}

\begin{enumerate}
  \item \textbf{车规IGBT/SiC放量}：V代IGBT+二代/四代SiC MOS适配800V平台，深度绑定主流车企，车载SiC累计出货超10万颗。
  \item \textbf{8英寸SiC产线}：士兰集宏（总投70亿），2025Q4通线，月产5000片，2026H2正式投产，面向8英寸SiC车规主驱。
  \item \textbf{12英寸高端模拟产线}：总投资200亿，2026年1月开工，2027Q4初步通线，远期年产24万片。
  \item \textbf{白电IPM持续领先}：国产替代空间+品类扩展（工业IPM、光伏IPM）。
  \item \textbf{功率半导体涨价周期}：行业景气度提升，产能利用率满载，议价能力增强。
  \item \textbf{全球份额提升}：从3.3\%向5\%+进军，产能+品类+车规三维驱动。
\end{enumerate}

\section{士兰微估值与市场定价}

\begin{tabularx}{\textwidth}{L{3cm}X}
\toprule
\textbf{Item} & \textbf{Value} \\
\midrule
最新股价 & 46.62元（2026-07-06），近期高位57.10元已回落 \\
总市值 & $\sim$776亿元 \\
PE TTM & $\sim$169x \\
PE（2026E） & 97--129x（基于6--8亿净利） \\
PE（2027E） & 50--65x（基于12--15亿净利） \\
机构目标价 & 凯基55元（+18\%）；证券之星相对估值38--42元 \\
近期走势 & 从高点57.10元回调至46.62元（-18\%），7/3主力净卖出5.67亿 \\
\bottomrule
\end{tabularx}

%% ================================================================
%% PART III: 对比分析
%% ================================================================
\newpage
\section*{\color{navy}\large Part III: 核心对比与投资结论}
\addcontentsline{toc}{section}{Part III: 对比与结论}

\section{维度对比}

\begin{tabularx}{\textwidth}{L{2.5cm}XX}
\toprule
\textbf{维度} & \textbf{华润微（688396）} & \textbf{士兰微（600460）} \\
\midrule
企业性质 & 央企（华润集团） & 民企 \\
市值 & $\sim$1152亿 & $\sim$776亿 \\
全球排名 & 功率半导体国内龙头 & 全球第六（3.3\%） \\
核心赛道 & MOSFET + AI电源 + 车规 & IGBT/IPM + 车规SiC + 白电 \\
最强单品 & MOSFET国内第一（>20\%） & 白电IPM全球第一（>40\%） \\
AI电源弹性 & 极高（全链路方案已批量） & 中等（12英寸模拟产线尚在建设） \\
SiC进度 & 650V--2300V已量产 & 6英寸已量产+8英寸2026H2投产 \\
产能节奏 & 重庆满产+深圳扩建 & 8英寸满载+8寸SiC通线+12寸开工 \\
2026E PE & 85--96x & 97--129x \\
业绩弹性 & 净利+297\%（Q1高基数效应后续减弱） & 净利+41\%（基数较高后增速稳健） \\
估值安全边际 & 中等（PE较低但已不便宜） & 偏低（估值较高+近期回调） \\
\bottomrule
\end{tabularx}

\section{Risks and Invalidation}

\begin{riskbox}[共同风险]
\begin{itemize}
  \item \textbf{行业周期风险}：功率半导体强周期属性，涨价持续性存疑。
  \item \textbf{估值透支}：两家公司PE均>80x（2026E），需持续高增验证。
  \item \textbf{海外竞争}：英飞凌、安森美在车规/AI领域仍占主导，国产替代非一蹴而就。
  \item \textbf{产能过剩}：国内功率半导体产能集中释放（2026--2027），可能引发价格战。
\end{itemize}
\end{riskbox}

\begin{riskbox}[华润微特有风险]
\begin{itemize}
  \item AI服务器电源从``批量供货''到``大规模放量''的转化节奏不确定。
  \item 12英寸参股项目盈利改善进度及资产注入预期存在不确定性。
  \item Q1高增主要因低基数，后续季度同比增速将自然放缓。
\end{itemize}
\end{riskbox}

\begin{riskbox}[士兰微特有风险]
\begin{itemize}
  \item 8英寸SiC产线2026H2才正式投产，上半年增量有限。
  \item 12英寸高端模拟产线需2027Q4才通线，短期无业绩贡献。
  \item 从高位57元急速回调至46元（-18\%），短期资金面承压。
  \item 机构目标价分歧大（36.5--55元），市场定价锚不稳。
\end{itemize}
\end{riskbox}

\section{Contrarian View}

\textbf{看空华润微}：央企IDM效率问题长期存在，毛利率25\%显著低于英飞凌（40\%+），即使营收增长但净利率提升缓慢。AI电源占比提升需要持续验证，若AI资本开支放缓则订单弹性消失。当前1150亿市值对标英飞凌市值（$\sim$5000亿人民币），隐含``中国英飞凌''故事能否兑现存疑。

\textbf{看空士兰微}：白电IPM市场天花板有限（$\sim$百亿级），SiC从6英寸到8英寸的良率爬坡存在不确定性。公司历史上多次``高增长+高估值+回调''循环，当前PE$\sim$169x处于历史高分位。12英寸200亿重资产投入期间，折旧压力将持续压制利润。

\textbf{核心验证指标}：
\begin{itemize}
  \item 华润微：Q2-Q3毛利率能否站稳30\%+，AI电源收入占比能否达30\%。
  \item 士兰微：8英寸SiC投产良率，2026H2车规SiC大批量出货确认。
  \item 共同：功率半导体涨价能否持续2+个季度，下游需求（新能源车/AI/光伏）是否超预期。
\end{itemize}

\section{投资结论与监控触发}

\subsection*{华润微（688396）}
\textbf{评级}：值得长期跟踪的核心资产，短期估值偏高但趋势向上。\\
\textbf{优选场景}：AI服务器电源国产替代主线、功率半导体涨价周期。\\
\textbf{关注价位}：回调至2026E PE$\sim$70x（对应股价$\sim$65--70元）时性价比改善。\\
\textbf{加仓信号}：Q2毛利率>30\%确认 + AI电源大客户新增订单公告。\\
\textbf{失效条件}：AI电源放量不及预期（Q3占比<25\%）或涨价周期1季度内结束。

\subsection*{士兰微（600460）}
\textbf{评级}：SiC+车规第二曲线具备高弹性，但短期回调未结束。\\
\textbf{优选场景}：新能源车800V SiC渗透率提升、功率半导体国产替代。\\
\textbf{关注价位}：回调至2027E PE$\sim$40x（对应股价$\sim$36--40元）时进入击球区。\\
\textbf{加仓信号}：8英寸SiC投产良率报告正面 + 车规SiC季度出货量翻倍确认。\\
\textbf{失效条件}：8英寸SiC良率低于预期（延期至2027）或车规大客户订单取消。

%% ================================================================
\vfill
\begin{disclosurebox}[Metadata]
\footnotesize
Confidence: 58/100 \quad | \quad
Data quality: partial source / web search only \quad | \quad
Degradation: Local data pipeline unavailable; analysis based on broker reports, public disclosures, and web aggregation.

\vspace{0.4em}
{\tiny\reportdisclaimer\ Generated by AI agents for research reference only. 2026-07-06.}
\end{disclosurebox}

\end{document}
